
\section{Proof-of-Concept Evaluation}

We instantiate LLM-GridEval in an EV load-altering attack scenario to demonstrate the framework's ability to compare static and adaptive attackers under identical conditions. The evaluation addresses two questions: (1) Can adaptive attackers exploit timing patterns to outperform random baselines? (2) What domain knowledge must LLM-based attackers possess to be effective?

\subsection{Experimental Setup}

\subsubsection{Grid Configuration}
The testbed uses the IEEE 123-node distribution feeder connected to an IEEE 9-bus transmission system via HELICS co-simulation. Six EV charging stations (EV1 through EV6) are distributed across phases A, B, and C at different nodes. The feeder operates with a base load of approximately 2.8 to 3.2 MW during the test period (simulation hour 7), with a protection threshold of $P_{\max} = 4.2$ MW. When exceeded, the defensive controller sheds EV load to restore safe operating conditions.

\subsubsection{Defensive Controller}
The threshold-based feeder controller polls total feeder power every 10 seconds. When $P_{\text{feeder}} \geq P_{\max}$, the controller commands all EV stations to zero power. When load returns below 2.6 MW, the controller gradually restores EV charging. This simple defense represents a baseline congestion management mechanism.

\subsubsection{Attack Constraints}
To create meaningful strategic trade-offs, we impose a scarce attack budget on all attacker variants:
\begin{itemize}
    \item \textbf{Cooldown}: Minimum 90 seconds between consecutive attacks (upper-bounding the rate to $\leq$40 attacks/hour)
    \item \textbf{Power bounds}: Each attack sets EV power within [500, 1500] kW
    \item \textbf{Ramping}: Power changes are rate-limited to 100 kW/s to prevent solver instability
\end{itemize}

These constraints ensure that each attack decision carries strategic weight: poor timing or target selection wastes scarce attack opportunities.

The 5-minute experiment window was selected to capture multiple attack-response cycles while keeping computational costs tractable for parameter sweeps. Each experiment begins at simulation hour 7, corresponding to moderate evening load conditions where the base load (2.8--3.2~MW) leaves headroom below the 4.2~MW threshold but is close enough that coordinated EV manipulation can trigger violations.

\subsection{Attack Variants}

We evaluate three attacker policies through the same MCP primitive interface, ensuring that differences in outcomes reflect policy quality rather than interface advantages.

\subsubsection{Random Baseline ($\pi_{\text{static}}$)}
The random attacker represents a static, non-adaptive policy. At each observation interval (5 seconds), if the cooldown has elapsed, the attacker attacks with 30\% probability. Target EV and power level are selected uniformly at random from the valid ranges. This baseline has no timing intelligence and no strategic planning.

\subsubsection{AI-V1: Timing-Only ($\pi_{\text{adapt}}^{(1)}$)}
The first adaptive policy uses an LLM with two-level timing intelligence:

\textbf{Macro-timing} assesses grid stress by computing the headroom between current load and threshold. A macro score of 100 indicates the threshold is already exceeded; scores above 70 indicate favorable attack conditions with headroom below 1 MW.

\textbf{Micro-timing} tracks the defender's observation cycle. The controller acts every 10 seconds; an attack issued immediately after a controller action (cycle position $\approx 0.0$) has nearly 10 seconds before the next defensive response, while an attack at cycle position 0.9 has only 1 second. The micro score maps cycle position to a 0 to 100 scale favoring early-cycle attacks.

The AI-V1 policy waits until the micro-timing score exceeds 70 before attacking, but receives no guidance on target selection or power accumulation dynamics.

\subsubsection{AI-V2: Timing with Strategy ($\pi_{\text{adapt}}^{(2)}$)}
The second adaptive policy extends V1 with explicit strategic knowledge. The LLM prompt includes:
\begin{itemize}
    \item \textbf{Power accumulation model}: Explanation that EV contributions are additive across different targets but overwrite when the same EV is attacked repeatedly
    \item \textbf{Diversification priority}: Instruction to attack unattacked EVs before repeating any target
    \item \textbf{Attack history}: Runtime context showing which EVs have been attacked and their current power levels
\end{itemize}

This additional context enables the LLM to reason about cumulative grid impact rather than optimizing individual attack timing in isolation.

\subsection{Power Accumulation Model}

A critical insight that emerged during V1 testing is that EV power contributions follow an additive-but-overwriting model:
\begin{equation}
P_{\text{total}} = P_{\text{base}} + \sum_{i=1}^{6} P_{\text{EV}_i}
\end{equation}

When an attacker issues a command to EV$_i$, it sets $P_{\text{EV}_i}$ to the specified value, overwriting any previous setpoint for that EV. Consequently, three attacks on EV1 at 1000, 1200, and 1500 kW contribute only 1500 kW total, while three attacks on EV1, EV2, and EV3 at 1500 kW each contribute 4500 kW. This property makes target diversification essential for effective load accumulation. We report the final accumulated EV setpoint $\sum_i P_{\text{EV}_i}$ at the end of each run; repeated attacks on the same EV overwrite its setpoint rather than add.

\subsection{Results}

Table~\ref{tab:results} summarizes results from 5-minute experiments under identical grid conditions and constraints. Each attacker variant was executed three times with different random seeds; we report mean values. The primary metric is Threshold Violation Duration (TVD): cumulative seconds during which feeder load exceeds $P_{\max}$. Higher TVD indicates more effective attacks from the adversary's perspective, and correspondingly, greater stress on the defensive controller. Attack success rate is the fraction of executed attacks after which $P_{\text{feeder}}$ exceeds $P_{\max}$ within the subsequent controller observation interval.

\begin{table}[t]
\centering
\caption{Comparison of Attack Strategies}
\label{tab:results}
\begin{tabular}{lccc}
\toprule
\textbf{Metric} & \textbf{Random} & \textbf{AI-V1} & \textbf{AI-V2} \\
\midrule
\multicolumn{4}{l}{\textit{Primary Metrics}} \\
TVD (seconds) & 240 & 120 & 240 \\
Total attacks executed & 4 & 3 & 3 \\
Attack success rate & 25\% & 33\% & 33\% \\
\midrule
\multicolumn{4}{l}{\textit{Timing Metrics}} \\
Avg. micro-timing score & 75 & 91 & 100 \\
Avg. cycle position & 0.25 & 0.09 & 0.00 \\
\midrule
\multicolumn{4}{l}{\textit{Strategic Metrics}} \\
Unique EVs targeted & 4 & 1 & 3 \\
Avg. power per attack (kW) & 865 & 1133 & 1500 \\
Final accumulated EV setpoint (kW) & 3459 & 1500 & 4500 \\
\bottomrule
\end{tabular}
\end{table}

\subsubsection{AI-V1 Analysis}
Despite achieving superior timing metrics (micro-score 91 vs. 75, cycle position 0.09 vs. 0.25), AI-V1 produced the worst TVD outcome at 120 seconds, half that of the random baseline. Examination of the attack log revealed that V1 attacked EV1 exclusively across all three attacks, with power levels of approximately 700, 1200, and 1500 kW. Each subsequent attack overwrote the previous setpoint, resulting in only 1500 kW of cumulative contribution despite expending three attacks.

This result demonstrates that timing optimization alone is insufficient. The LLM optimized a local objective (attack at good timing) without understanding the global objective (maximize cumulative load). From the LLM's perspective, EV1 was always a valid target, and attacking it at optimal timing appeared to be a good decision.

\subsubsection{AI-V2 Analysis}
With explicit strategic guidance, AI-V2 achieved both perfect timing (micro-score 100, cycle position 0.00) and effective diversification. The attack sequence targeted EV1, then EV2, then EV3, each at maximum power (1500 kW). This produced 4500 kW of cumulative EV load, significantly exceeding the random baseline's 3459 kW.

AI-V2 matched random's TVD of 240 seconds while using 25\% fewer attacks (3 vs. 4). The per-attack efficiency of 80 seconds TVD per attack for V2 versus 60 seconds for random reflects V2's superior resource utilization. Under the scarce budget constraints, this efficiency advantage would compound over longer attack campaigns.

\subsubsection{Ceiling Effect}
Both random and AI-V2 achieved identical TVD despite V2's superior metrics. This ceiling effect arises because the grid at Hour 7 is sufficiently stressed that even suboptimal attacks can trigger violations. Under harder scenarios, such as higher thresholds, lower base load, or faster controller response, the evaluation validity gap $\Delta E(M)$ between adaptive and static attackers would widen.

\subsection{Discussion}
These experiments yield three insights for LLM-based cyber-physical attackers. First, tactical excellence does not guarantee strategic success: AI-V1's timing optimization represented locally optimal decisions that were globally suboptimal, a failure mode characteristic of LLMs operating without sufficient domain context. This suggests that prompt engineering for cyber-physical attacks must go beyond protocol-level instructions to include system dynamics.

Second, domain model awareness is essential: AI-V2 succeeded only after receiving explicit knowledge of power accumulation dynamics, demonstrating that cyber-physical LLM attackers require grounding in physical models, not just protocol knowledge. For defenders, this implies that attacks requiring multi-step physical reasoning may be detectable through their characteristic patterns of information gathering.

Third, the framework enables systematic attacker development: the V1-to-V2 evolution illustrates how LLM-GridEval supports iterative refinement of attacker policies within a controlled environment, with each variant using identical primitives and constraints to isolate the impact of policy changes. This capability is essential for rigorous security evaluation where confounding variables must be minimized.

The current experiments represent initial validation under favorable attack conditions. Future work will evaluate longer campaigns (1--2 hours), harder grid configurations, and additional attack primitives to more fully characterize the evaluation validity gap.
